\section{Návrh riešenia}
\label{sec:navrh}

% ----- Požiadavky -----
\subsection{Požiadavky}

Systém musí spĺňať tieto požiadavky:
% Odrážkový zoznam.
\begin{itemize}
  \item zaznamenať pre každé stlačenie klávesy čas stlačenia aj uvoľnenia
        s rozlíšením lepším ako milisekunda,
  \item zabezpečiť, aby všetci účastníci písali tú istú vetu rovnakým
        spôsobom (bez opráv a bez vkladania textu),
  \item uložiť namerané údaje v tvare vhodnom na ďalšie spracovanie,
  \item z jedného nového pokusu určiť najpodobnejších účastníkov a posúdiť,
        či ide o neznámu osobu,
  \item umožniť vyhodnotenie úspešnosti na dátach, ktoré neboli použité
        na tvorbu šablón.
\end{itemize}

% ----- Architektúra -----
\subsection{Architektúra}

Riešenie má dve časti (Obr.~\ref{fig:arch}). Webová aplikácia
(klient) beží v prehliadači, zaznamenáva klávesové udalosti a zobrazuje
výsledky. Serverová časť (REST API) ukladá relácie do databázy SQLite
a obsahuje algoritmus identifikácie. Komunikácia prebieha vo formáte JSON
cez HTTP.

% Obrázok (schéma nakreslená v TikZ): Architektúra systému.
\begin{figure}[ht]
  \centering
  \vspace{1.5em}   % <- prázdny riadok (medzera) medzi textom hore a schémou; zmeňte číslo podľa potreby
  \begin{tikzpicture}[
      box/.style={draw, rounded corners=4pt, thick, minimum width=3.4cm, minimum height=1.9cm,
                  align=center, fill=gray!10},
      lbl/.style={font=\small, align=center, inner sep=2pt},
      cap/.style={font=\footnotesize\itshape, text=gray},
      >={Stealth[length=3mm]}, node distance=2.6cm]
    % tri bloky vedľa seba: klient -> server -> databáza
    \node[box] (ui)  {\textbf{Webová aplikácia}\\[3pt]\small React, prehliadač};
    \node[box, right=of ui]  (api) {\textbf{REST API}\\[3pt]\small FastAPI, Python};
    \node[box, right=of api] (db)  {\textbf{Databáza}\\[3pt]\small SQLite};
    % šípky v oboch smeroch (požiadavka a odpoveď)
    \draw[<->, thick] (ui) -- node[above, lbl] {JSON\\cez HTTP} (api);
    \draw[<->, thick] (api) -- node[above, lbl] {SQL} (db);
    % popisky pod blokmi
    \node[cap, below=4pt of ui]  {klient};
    \node[cap, below=4pt of api] {server};
    \node[cap, below=4pt of db]  {úložisko};
  \end{tikzpicture}
  \caption{Architektúra systému.}
  \label{fig:arch}
\end{figure}

Klient pracuje v štyroch stavoch: príprava relácie, meranie, identifikácia
jedného pokusu a záverečná obrazovka s možnosťou odoslať dáta na server.

% ----- Zber údajov -----
\subsection{Zber údajov}

Účastník zadá meno, vyberie zariadenie (notebook alebo stolný počítač)
a typ klávesnice (vstavaná, externá USB, externá Bluetooth) a vyjadrí súhlas
so záznamom. Následne \textbf{15-krát} napíše vetu
\emph{„Every keystroke tells a unique story.“} (37 znakov) a každý pokus
potvrdí klávesom Enter.

Na meranie času sa používa \code{performance.now()}. Čas každej udalosti sa
ukladá relatívne k prvému stlačeniu v danom pokuse. Klávesy Shift, Caps Lock,
Ctrl, Alt, Meta, Tab a Escape sa nezaznamenávajú. Opakované udalosti pri
podržaní klávesy sa blokujú. Vkladanie, vystrihovanie, presúvanie textu
a kontextové menu sú vo vstupnom poli zakázané.

Pokus sa zamietne a účastník ho musí zopakovať, ak:
% Odrážkový zoznam.
\begin{itemize}
  \item napísaný text sa nezhoduje s vetou,
  \item počet stlačení sa líši od počtu uvoľnení (nie všetky klávesy boli
        uvoľnené).
\end{itemize}
Na strane servera sa navyše pri vytváraní šablón ignoruje každý pokus,
ktorý nemá presne 37 stlačení a 37 uvoľnení (podkapitola~\ref{sec:priznaky}).

% ----- Formát údajov -----
\subsection{Formát údajov}

Každá klávesová udalosť je jeden riadok s poľami uvedenými v
Tab.~\ref{tab:format}. Relácia sa na server odosiela ako objekt
\code{\{meta, rows\}}, kde \code{meta} obsahuje identifikáciu účastníka
a relácie a \code{rows} zoznam udalostí všetkých pokusov.

% Tabuľka: Polia záznamu o klávesovej udalosti.
\begin{table}[ht]
  \centering
  \small
  \caption{Polia záznamu o klávesovej udalosti.}
  \label{tab:format}
  \begin{tabularx}{\textwidth}{|l|Y|}
    \hline
    \rowcolor[gray]{0.9} \textbf{Pole} & \textbf{Význam} \\
    \hline
    \code{participant\_id} & identifikátor účastníka (zadané meno) \\
    \hline
    \code{session\_id} & čas začiatku relácie vo formáte ISO 8601 \\
    \hline
    \code{device}, \code{keyboard} & zariadenie a typ klávesnice \\
    \hline
    \code{rep} & poradové číslo pokusu (1--15) \\
    \hline
    \code{seq} & poradie udalosti v rámci pokusu \\
    \hline
    \code{event} & \code{down} (stlačenie) alebo \code{up} (uvoľnenie) \\
    \hline
    \code{code}, \code{key} & fyzický kód klávesy a vyprodukovaný znak \\
    \hline
    \code{t\_ms} & čas v ms od prvého stlačenia v pokuse \\
    \hline
  \end{tabularx}
\end{table}

% ----- Príznaky -----
\subsection{Príznaky}
\label{sec:priznaky}

Z jedného pokusu s $n=37$ klávesmi sa vytvorí vektor príznakov. Nech $d_i$ je
čas $i$-teho stlačenia a $u_i$ čas uvoľnenia tej istej klávesy (prvá udalosť
\code{up} s rovnakým \code{code}, ktorá nasleduje po stlačení). Potom platia
vzorce~\eqref{eq:hold} a~\eqref{eq:gap}:
% Vzorec (v texte sa číta ako matematický zápis).
\begin{align}
  h_i &= u_i - d_i, & i &= 1,\dots,n, \label{eq:hold}\\
  g_i &= d_{i+1} - d_i, & i &= 1,\dots,n-1, \label{eq:gap}
\end{align}
kde $h_i$ je doba držania klávesy (vzorec~\eqref{eq:hold}) a $g_i$ interval medzi
dvoma po sebe idúcimi stlačeniami (vzorec~\eqref{eq:gap}). Vektor príznakov
$\mathbf{x} = (h_1,\dots,h_n,g_1,\dots,g_{n-1})$
má $m = 2n-1 = 73$ zložiek. Pokus s iným počtom stlačení či uvoľnení,
s chýbajúcim uvoľnením, nečíselným alebo záporným časom je neplatný.

% ----- Šablóna účastníka a miera podobnosti -----
\subsection{Šablóna účastníka a miera podobnosti}

Šablóna účastníka $P$ je dvojica vektorov $(\boldsymbol{\mu}, \boldsymbol{\sigma})$,
ktoré sa vypočítajú ako priemer a smerodajná odchýlka (populačná,
$\mathrm{ddof}=0$) jednotlivých príznakov z pokusov č.~1 až 10 všetkých
relácií daného účastníka. Ak je odchýlka príznaku nulová, nahradí sa
hodnotou 1, aby sa predišlo deleniu nulou.

Vzdialenosť nového pokusu $\mathbf{x}$ od šablóny $P$ je štandardizovaná
euklidovská vzdialenosť
% Vzorec (v texte sa číta ako matematický zápis).
\begin{equation}
  \label{eq:dist}
  d(\mathbf{x}, P) = \sqrt{\sum_{j=1}^{m}\left(\frac{x_j-\mu_j}{\sigma_j}\right)^{2}}.
\end{equation}
Vo vzorci~\eqref{eq:dist} je $m$ počet príznakov a $\mu_j$, $\sigma_j$ sú priemer
a smerodajná odchýlka $j$-teho príznaka v šablóne.
Účastníci sa zoradia podľa vzdialenosti vzostupne. Pokus sa označí ako
\emph{neznámy}, ak pre najbližšieho účastníka platí
% Vzorec (v texte sa číta ako matematický zápis).
\begin{equation}
  \label{eq:unknown}
  \frac{d_{\min}}{\sqrt{m}} > 3,
\end{equation}
Podmienka~\eqref{eq:unknown} znamená, že priemerná štandardizovaná odchýlka na jeden
príznak je väčšia ako tri.
Prah 3 je zvolený empiricky a nebol kalibrovaný na dátach.

Pre troch najbližších účastníkov sa vypočíta relatívna miera podobnosti
% Vzorec (v texte sa číta ako matematický zápis).
\begin{equation}
  \label{eq:share}
  p_i = 100\cdot\frac{w_i}{\sum_{k=1}^{3} w_k}, \qquad w_i = \frac{1}{d_i + 10^{-9}}.
\end{equation}
Hodnota $p_i$ zo vzorca~\eqref{eq:share} vyjadruje len relatívne poradie medzi tromi najlepšími
kandidátmi, \textbf{nie pravdepodobnosť}, že ide o danú osobu.

% ----- Číselný príklad výpočtu -----
\subsection{Číselný príklad výpočtu}
\label{sec:priklad}

Postup z predchádzajúcich podkapitol je ilustrovaný na zjednodušenom príklade
s tromi klávesmi. V skutočnom pokuse je $n=37$ a vektor má 73 zložiek,
tu je $n=3$ a vektor má $m=2n-1=5$ zložiek. Všetky hodnoty sú vymyslené
a slúžia len na vysvetlenie.

Pokus obsahuje udalosti z Tab.~\ref{tab:priklad-udalosti}.

% Tabuľka: Ilustračný pokus s tromi klávesmi (časy v ms).
\begin{table}[ht]
  \centering
  \small
  \caption{Ilustračný pokus s tromi klávesmi (časy v ms).}
  \label{tab:priklad-udalosti}
  \begin{tabularx}{\textwidth}{|C|C|C|C|}
    \hline
    \rowcolor[gray]{0.9} \textbf{Klávesa} & \textbf{Stlačenie $d_i$} & \textbf{Uvoľnenie $u_i$} & \textbf{Držanie $h_i=u_i-d_i$} \\
    \hline
    1 & 0 & 95 & 95 \\
    \hline
    2 & 140 & 230 & 90 \\
    \hline
    3 & 310 & 385 & 75 \\
    \hline
  \end{tabularx}
\end{table}

Intervaly medzi stlačeniami sú $g_1 = 140-0 = 140$ a $g_2 = 310-140 = 170$,
teda $\mathbf{x} = (95, 90, 75, 140, 170)$.

Nech šablóna účastníka obsahuje priemery a smerodajné odchýlky
$\boldsymbol{\mu} = (100, 90, 80, 150, 160)$ a
$\boldsymbol{\sigma} = (10, 10, 5, 20, 20)$. Štandardizované odchýlky sú
% Vzorec (v texte sa číta ako matematický zápis).
\[
  z_j = \frac{x_j-\mu_j}{\sigma_j} = (-0{,}5;\ 0;\ -1;\ -0{,}5;\ 0{,}5),
\]
a po dosadení do vzorca~\eqref{eq:dist} vzdialenosť
% Vzorec (v texte sa číta ako matematický zápis).
\[
  d = \sqrt{0{,}25 + 0 + 1 + 0{,}25 + 0{,}25} = \sqrt{1{,}75} \approx 1{,}323 .
\]
Pre kontrolu neznámeho pokusu podľa vzorca~\eqref{eq:unknown} je
$d/\sqrt{m} = 1{,}323/\sqrt{5} \approx 0{,}592$, čo je výrazne menej ako
prah 3, takže pokus sa neoznačí ako neznámy.

Ak by boli vzdialenosti od šablón troch účastníkov $d_1 \approx 1{,}323$,
$d_2 = 2{,}0$ a $d_3 = 4{,}0$, váhy sú $w_i = 1/d_i$ približne
$0{,}756$; $0{,}5$ a $0{,}25$ a relatívne miery podobnosti podľa vzorca~\eqref{eq:share}
$p_1 \approx 50{,}2\,\%$, $p_2 \approx 33{,}2\,\%$ a $p_3 \approx 16{,}6\,\%$.
Najbližší účastník má teda najvyššie percento, no hodnota 50,2\,\% neznamená,
že pravdepodobnosť zhody je 50,2\,\%.

% ----- Osobné údaje a etika zberu -----
\subsection{Osobné údaje a etika zberu}
\label{sec:etika}

Dynamika písania je biometrický údaj, ktorý môže jednoznačne súvisieť
s konkrétnou osobou, a preto sa so zozbieranými dátami zaobchádza ako
s osobnými údajmi. V systéme sa uplatňujú tieto zásady:
% Odrážkový zoznam.
\begin{itemize}
  \item \textbf{Informovaný súhlas.} Meranie sa nedá spustiť bez zaškrtnutia
        súhlasu so záznamom časových údajov stlačenia kláves na výskumné
        účely.
  \item \textbf{Minimalizácia údajov.} Zaznamenávajú sa len časy stlačenia
        a uvoľnenia kláves jednej vopred určenej vety. Obsah písaného textu
        je pre všetkých rovnaký a neobsahuje súkromné informácie. Okrem
        mena (identifikátora účastníka) sa ukladá len typ zariadenia
        a klávesnice.
  \item \textbf{Identifikátor účastníka.} Aplikácia vyžaduje meno
        a priezvisko, ktoré sa používa ako \code{participant\_id}. Databáza
        teda obsahuje osobné údaje a nemá sa zverejňovať.
  \item \textbf{Lokálne úložisko.} Dáta sa ukladajú do súboru SQLite
        na serveri riešenia; nepoužívajú sa žiadne externé služby.
  \item \textbf{Zdieľanie kódu bez dát.} Do repozitára patrí kód, nie
        databáza.
\end{itemize}

Na zverejnenie alebo odovzdanie dát tretím osobám treba zvážiť, či účastníci
s takým použitím súhlasili. Tento dokument nepredstavuje právne posúdenie.

% ----- Protokol vyhodnotenia -----
\subsection{Protokol vyhodnotenia}
\label{sec:protokol}

Aby sa šablóny neotestovali na tých istých dátach, z ktorých vznikli, sa
pokusy rozdelia podľa poradia: pokusy 1--10 slúžia na tvorbu šablóny,
pokusy 11--15 na testovanie. Pre každý testovací pokus sa zistí poradie
účastníkov a vyhodnotia sa:
% Odrážkový zoznam.
\begin{itemize}
  \item \textbf{top-1}: podiel pokusov, pri ktorých je skutočný účastník
        na prvom mieste (pokus označený ako neznámy sa počíta ako chybný),
  \item \textbf{top-3}: podiel pokusov, pri ktorých je skutočný účastník
        medzi tromi najpodobnejšími,
  \item \textbf{matica zámen}: počty predikcií pre každého skutočného
        účastníka.
\end{itemize}
